\section{Background}\label{sec:background}
\subsection{Industrial Control Systems and PLCs}
Modern ICS architectures include Supervisory Control and Data Acquisition (SCADA) systems, Distributed Control Systems (DCS), Programmable Logic Controllers (PLCs), Remote Terminal Units (RTUs), Human Machine Interfaces (HMIs), and industrial networks that often keep legacy protocols designed without security mechanisms \cite{Zhu2011Taxonomy}. Their integration with Ethernet and Internet-based technologies improved interoperability but increased cybersecurity exposure \cite{Bhamare2020ICS,Kayan2021ICPS}. Industry 4.0 and IIoT paradigms connect industrial devices through standard TCP/IP protocols and cloud services \cite{Souza2021SmartManufacturing}, which expands the attack surface to malware, unauthorized access, and denial-of-service attacks \cite{Pedreira2021Industry40}.

PLCs automate electromechanical processes in real time: they read sensor inputs, run the control logic, and drive actuators in manufacturing, energy, and process control \cite{Makrakis2021ICS}. Many PLCs were designed for isolated environments and present weak authentication, insecure firmware updates, unencrypted communications, and insufficient traffic filtering, which exposes the process to command injection, unauthorized access, and denial of service \cite{Nicholson2021ICSAttacks}.

\subsection{Denial-of-Service (DoS) Attacks}
Common DoS techniques in industrial networks are SYN, ICMP, and UDP floods and protocol-specific floods against Modbus/TCP. They fall into bandwidth exhaustion, resource exhaustion, or protocol vulnerability exploitation \cite{Specht2004DDoS}, and may exhaust CPU, saturate links, or overload buffers. Incident reports describe flooding of the communication stack of widely used PLCs that causes restarts or loss of process control \cite{Chabrow2016ICS}. The effect of a flood can be summarized as $R_{attack} > R_{processing}$: the traffic generation rate exceeds the processing capacity of the device or of the network element that has to forward it.

\subsection{Intrusion Prevention Systems (IPS)}
Unlike intrusion detection systems, an IPS works inline in the communication path and can inspect, filter, and mitigate in real time \cite{Raman2021ICS}. Industrial IPS solutions combine signature-based detection, anomaly-based detection, behavioral analysis, and protocol-aware inspection \cite{Yu2019IDS}, which are the standard classification of detection methods \cite{Scarfone2007IDS}. Signature-based approaches match known patterns with rules; anomaly-based approaches use statistical or learning models, and keeping false positives low usually requires reference models of the industrial protocol \cite{Liao2013IDS}.

Industrial protocols such as Modbus/TCP, EtherNet/IP, Profinet, DNP3, and OPC UA carry the exchange between PLCs, SCADA, sensors, and actuators \cite{Pan2020PLC}. Modbus/TCP is widely used for its simplicity but has no authentication, encryption, or integrity checks, so it is exposed to interception, replay, command injection, and denial of service \cite{Yu2019IDS}. Industrial security strategies therefore combine IPS, firewalls, and VLAN segmentation \cite{Lin2021Protocols}.

\subsection{Inline Prevention with Suricata and NFQUEUE}
An IPS can be deployed inline on a Linux router by diverting forwarded packets from the kernel packet filter to a user-space engine. In this work the \texttt{nftables} \texttt{queue} statement sends the forwarded packets to \texttt{NFQUEUE}, and Suricata returns a verdict for each packet: \emph{accept}, or \emph{drop} when a rule with the \texttt{drop} action matches. With the \texttt{alert} action Suricata only logs the match and accepts the packet, so any packet that does not reach its destination in that mode was lost for another reason, such as an overflowing queue. An inline IPS therefore adds an element to the data path whose capacity bounds the traffic the network can forward \cite{fraunholz2018industrial}.

\subsection{Related Work and Positioning}
Surveys of intrusion detection in industrial control systems \cite{Liao2013IDS,Bhamare2020ICS,ahmed2019survey} describe signature-based, anomaly-based, and protocol-aware techniques, and most experimental studies evaluate detection models, often based on machine learning \cite{Tama2022IDS,Raman2021ICS}. Fraunholz et al.\ propose and implement a concept for reactive protection of industrial networks with an IPS \cite{fraunholz2018industrial}, and Yu et al.\ propose a two-stage intrusion detection system for Ethernet/IP networks \cite{Yu2019IDS}. In contrast with detection-centred studies, this work measures what an inline, rule-based IPS does to the availability of a PLC service, using a reproducible free-software testbed, a baseline that isolates the firewall and the segmentation, repeated runs with statistical analysis, and independent sources for the traffic that the IPS removes.
